% © 2026 Ing. David Jaroš — CC BY-NC-ND 4.0
%
% This work is licensed under a Creative Commons Attribution-NonCommercial-NoDerivatives
% 4.0 International License (CC BY-NC-ND 4.0).
%
% Spectral 3/2 Program — Step 2: Θ-Operator Candidates
% Track: SPECTRAL — independent of Hecke/twin-prime route
% Date: 2026-03-10

\ifdefined\INCLUDEMODE
  % Being included in another document — skip preamble
\else
  \documentclass[11pt,a4paper]{article}
  \usepackage{amsmath,amssymb,amsthm}
  \usepackage{hyperref}
  \usepackage{booktabs}

  \newtheorem{theorem}{Theorem}
  \newtheorem{lemma}{Lemma}
  \newtheorem{proposition}{Proposition}
  \newtheorem{conjecture}{Conjecture}
  \theoremstyle{definition}
  \newtheorem{definition}{Definition}
  \newtheorem{gap}{Gap}
  \theoremstyle{remark}
  \newtheorem{remark}{Remark}

  \title{Θ-Operator Candidates for the Spectral 3/2 Program\\
    \large Identifying the Natural Laplace Operator on the Θ-Field Configuration Space}
  \author{Ing.\ David Jaroš}
  \date{2026-03-10}

  \begin{document}
  \maketitle
\fi

% ======================================================================
\section{Introduction and Scope}
\label{sec:intro}
% ======================================================================

\textbf{Program context.}
This document is Step~2 of the Spectral 3/2 Program, which investigates
whether the exponent $3/2$ in
\begin{equation}
  B_{\mathrm{base}} = N_{\mathrm{eff}}^{3/2} \approx 41.57
  \label{eq:B_conj}
\end{equation}
arises from the spectral geometry of the Θ-field in UBT.
The exponent $3/2$ is currently a \textbf{[MOTIVATED CONJECTURE]};
see \texttt{consolidation\_project/alpha\_derivation/b\_base\_hausdorff.tex}.

\textbf{Independence statement.}
This track does not use $\alpha^{-1} = 137$ as an input.
The only inputs are: $N_{\mathrm{eff}} = 12$ (\textbf{[PROVED]}),
$\dim_\mathbb{R}(\mathrm{Im}\,\mathbb{H}) = 3$ (\textbf{[PROVED]}),
and the UBT action $\mathcal{S}[\Theta]$.

\textbf{Goal of this document.}
Identify the most natural Laplace-type or Hamiltonian operator associated
with the Θ-field, and determine whether it factorises into independent
spacetime and internal parts.

% ======================================================================
\section{The UBT Action and Kinetic Term}
\label{sec:action}
% ======================================================================

The UBT action functional is (see \texttt{THETA\_FIELD\_DEFINITION.md §5.1}):
\begin{equation}
  \mathcal{S}[\Theta]
  \;=\; \int_{\mathcal{M}} d^4q\, d^2\tau\, \sqrt{|\!\det G|}
    \Bigl[
      \tfrac{1}{2}\bigl\langle \nabla_\mu \Theta,\, \nabla^\mu \Theta \bigr\rangle
      - V(\Theta)
    \Bigr],
  \label{eq:S_full}
\end{equation}
where the integration domain is $\mathcal{M} = \mathbb{R}^4 \times S^1_\psi$
($S^1_\psi$: the compact imaginary-time circle of radius $R_\psi$ with
periodicity $\psi \sim \psi + 2\pi$).

Expanding around a background $\Theta_0 + \delta\Theta$, the quadratic
fluctuation action is:
\begin{equation}
  \mathcal{S}^{(2)}[\delta\Theta]
  \;=\; \tfrac{1}{2} \int d^4q\, d^2\tau
    \bigl\langle \delta\Theta,\, \hat{\mathcal{O}}\,\delta\Theta \bigr\rangle,
  \label{eq:S_quad}
\end{equation}
where the \emph{fluctuation operator} is:
\begin{equation}
  \hat{\mathcal{O}} \;:=\; -\nabla^\dagger \nabla + \mathcal{V}''(\Theta_0),
  \label{eq:O_def}
\end{equation}
with $\mathcal{V}''(\Theta_0)$ denoting the second variation of the potential.
In the free-field or vacuum-polarisation limit, $\mathcal{V}'' \approx m^2$
(mass term), so $\hat{\mathcal{O}} \approx -\nabla^\dagger\nabla + m^2$.

\begin{definition}[Θ-Laplacian]
\label{def:Delta_Theta}
  The \emph{Θ-Laplacian} is the principal part of $\hat{\mathcal{O}}$:
  \begin{equation}
    \Delta_\Theta \;:=\; \nabla^\dagger\nabla
      \;=\; G^{\mu\nu} \nabla_\mu \nabla_\nu,
    \label{eq:Delta_Theta}
  \end{equation}
  acting on sections of the field bundle
  $\mathcal{E} = \mathcal{E}_{\mathrm{biquat}} \otimes \mathcal{E}_{\mathrm{spinor}}
   \otimes \mathcal{E}_{\mathrm{gauge}}$.
  The fluctuation operator is $\hat{\mathcal{O}} = -\Delta_\Theta + m^2$.
\end{definition}

% ======================================================================
\section{Decomposition of the Domain}
\label{sec:domain}
% ======================================================================

The domain of $\Delta_\Theta$ decomposes as:
\begin{equation}
  \mathcal{M} \;=\; \underbrace{\mathbb{R}^4}_{\text{spacetime}}
    \;\times\; \underbrace{S^1_\psi}_{\text{compact }\psi}
  \label{eq:M_decomp}
\end{equation}
and the field bundle decomposes along the quaternion structure as:
\begin{equation}
  \delta\Theta \;=\; \underbrace{\delta\Theta_S}_{\in\,\mathbb{R}}
    \;\oplus\; \underbrace{\delta\Theta_V}_{\in\,\mathrm{Im}(\mathbb{H})}
  \label{eq:field_decomp}
\end{equation}
where $\mathrm{Im}(\mathbb{H}) = \mathrm{span}_\mathbb{R}\{I, J, K\}$ is the
pure-imaginary quaternion subspace (dimension 3, \textbf{[PROVED]}).

By gauge invariance the vacuum polarisation tensor is transverse,
$\Pi_{\mu\nu} \propto (g_{\mu\nu} - k_\mu k_\nu/k^2)$, and
$\delta\Theta_S$ (the scalar part) is a singlet under the gauge group and
decouples from $\Pi_{\mu\nu}$.
\textbf{Only $\delta\Theta_V$ contributes to the gauge coupling running.}
(Proved; see \texttt{b\_base\_hausdorff.tex} Lemma~1.)

% ======================================================================
\section{Three Candidate Operators}
\label{sec:candidates}
% ======================================================================

\subsection{Candidate 1: Full Spacetime Laplacian $\Delta^{(4+1)}$}
\label{sec:cand1}

The most na\"ive choice is the Laplacian on the full domain
$\mathbb{R}^4 \times S^1_\psi$:
\begin{equation}
  \Delta^{(4+1)} \;=\; \partial_t^2 + \partial_{x_1}^2 + \partial_{x_2}^2
    + \partial_{x_3}^2 + \partial_\psi^2.
  \label{eq:Delta_41}
\end{equation}
This operator has effective spectral dimension $d_{\mathrm{eff}} = 4$
on $\mathbb{R}^4$ (or $d_{\mathrm{eff}} = 5$ on the full product space
including $S^1_\psi$).

\textbf{Heat kernel:} $K^{(4+1)}(t) \sim (4\pi t)^{-5/2}$ for small $t$
on $\mathbb{R}^4 \times S^1_\psi$, giving the Weyl count
$N^{(4+1)}(\lambda) \sim C\,\lambda^{5/2}$.

\textbf{Assessment for B\_base:} The exponent $5/2 \neq 3/2$.
This candidate does not produce the correct scaling.

\subsection{Candidate 2: Compact-$\psi$ Laplacian $\Delta^{(S^1)}$}
\label{sec:cand2}

If one integrates over the non-compact $\mathbb{R}^4$ directions and retains
only the compact $\psi$-Kaluza–Klein tower, the relevant operator is:
\begin{equation}
  \Delta^{(S^1)} \;=\; -\partial_\psi^2,
  \qquad \psi \sim \psi + 2\pi,
  \label{eq:Delta_S1}
\end{equation}
with eigenvalues $\lambda_n = n^2/R_\psi^2$ ($n \in \mathbb{Z}$).

\textbf{Heat kernel:} $K^{(S^1)}(t) \sim (4\pi t)^{-1/2}$ (1-dimensional),
giving $N^{(S^1)}(\lambda) \sim C\,\lambda^{1/2}$.

\textbf{Assessment for B\_base:} The exponent $1/2 \neq 3/2$.
This candidate also fails.
Moreover, the Kaluza–Klein sum approach (approach A1) was already found to
diverge without a natural cutoff; see
\texttt{consolidation\_project/alpha\_derivation/b\_base\_spinor\_approach.tex}.

\subsection{Candidate 3: Internal Im($\mathbb{H}$) Laplacian $\Delta^{(3)}$}
\label{sec:cand3}

After the scalar decoupling ($\delta\Theta_S$ irrelevant) and
restricting to the vacuum-polarisation sector, the relevant fluctuations
$\delta\Theta_V$ live in the \emph{internal} space
$\mathrm{Im}(\mathbb{H}) \cong \mathbb{R}^3$.

If one treats $\mathrm{Im}(\mathbb{H})$ as carrying the effective
mode-counting dimension, the natural operator is the Laplacian on
a 3-dimensional mode space:
\begin{equation}
  \Delta^{(3)}_{\mathrm{Im}\,\mathbb{H}}
  \;=\; \partial_{I}^2 + \partial_{J}^2 + \partial_{K}^2,
  \label{eq:Delta_3}
\end{equation}
where $\partial_I, \partial_J, \partial_K$ are derivatives in the
$I, J, K$ quaternion-imaginary directions.

\textbf{Heat kernel:} $K^{(3)}(t) \sim (4\pi t)^{-3/2}$,
giving $N^{(3)}(\lambda) \sim C\,\lambda^{3/2}$.

\textbf{Assessment for B\_base:} The exponent $3/2$ is \emph{exactly}
what is needed.
\textbf{This is the primary candidate for the spectral 3/2 mechanism.}

% ======================================================================
\section{Factorisation Test}
\label{sec:factorisation}
% ======================================================================

The full Θ-Laplacian on $\mathbb{R}^4 \times S^1_\psi$ acts on
both spacetime and internal indices.
We test whether it factorises as:
\begin{equation}
  \Delta_\Theta \;=\; \Delta_{\mathbb{R}^4} \,\otimes\, \mathbf{1}_{\mathrm{int}}
    \;+\; \mathbf{1}_{\mathbb{R}^4} \,\otimes\, \Delta_{\mathrm{int}}
  \label{eq:factorisation}
\end{equation}
where $\Delta_{\mathrm{int}}$ acts on the internal (gauge/quaternion) indices.

\begin{proposition}[Factorisation in flat space — CANDIDATE]
\label{prop:factorise}
  In flat Minkowski space $\mathbb{R}^{3,1}$ with no background field
  ($\Theta_0 = \mathrm{const}$, $V'' = m^2$), the fluctuation operator
  factorises as:
  \begin{equation}
    \hat{\mathcal{O}}
    \;=\; \bigl(-\partial_\mu\partial^\mu + m^2\bigr)
      \,\otimes\,
      \mathbf{1}_{\mathrm{Im}\,\mathbb{H}},
    \label{eq:O_flat}
  \end{equation}
  so that each of the three $\mathrm{Im}(\mathbb{H})$-components obeys
  the same scalar Klein–Gordon equation.

  \begin{proof}[Proof sketch]
    In flat space, $G^{\mu\nu} = \eta^{\mu\nu}$, and
    $\nabla_\mu = \partial_\mu$ (no connection).
    The potential $V(\Theta) = (m^2/2)|\Theta_V|^2$ is diagonal in the
    $\{I,J,K\}$ basis by the $\mathrm{SO}(3)$ symmetry of $\mathrm{Im}(\mathbb{H})$.
    Hence $\hat{\mathcal{O}}$ is proportional to the identity on
    $\mathrm{Im}(\mathbb{H})$.
    The spectrum of each component is that of a scalar field on $\mathbb{R}^4$;
    the multiplicity from $\mathrm{Im}(\mathbb{H})$ is exactly 3.
  \end{proof}
\end{proposition}

\begin{remark}[Status of factorisation]
  The factorisation~\eqref{eq:O_flat} holds in flat space and is
  \textbf{[PROVED AS A FLAT-SPACE IDENTITY]}.
  In curved space (non-trivial $G_{\mu\nu}$), mixing between spacetime and
  internal directions may arise from spin connections; this case is not
  treated here.
\end{remark}

% ======================================================================
\section{Role of the Compact $\psi$-Direction}
\label{sec:compact_psi}
% ======================================================================

The compact imaginary-time direction $\psi \sim \psi + 2\pi$ contributes
a Kaluza–Klein tower to the spectrum.
Its contribution to the heat kernel is:
\begin{equation}
  K^{(S^1)}(t) \;=\; \sum_{n \in \mathbb{Z}} e^{-t\,n^2/R_\psi^2}
  \;\sim\; \frac{\sqrt{\pi}\,R_\psi}{\sqrt{t}}, \quad t \to 0^+,
  \label{eq:K_S1}
\end{equation}
by the Poisson summation formula (see also
\texttt{ubt\_with\_chronofactor/scripts/spectral/poisson\_duality\_demo.py}).

If the full domain includes the $S^1_\psi$ factor, the combined small-$t$
asymptotics of $\Delta^{(4+1)}$ give $K \sim t^{-5/2}$ (5 effective
dimensions), \emph{not} $t^{-3/2}$.

\textbf{Conclusion for this section:}
The $\psi$-direction does not produce the desired $t^{-3/2}$ on its own.
For the spectral $3/2$ mechanism, the relevant manifold must be 3-dimensional,
not 1-dimensional.
This supports Candidate~3 (Section~\ref{sec:cand3}) as the operative one,
provided the mode counting is governed by the internal $\mathrm{Im}(\mathbb{H})$
structure and not by the $\psi$-compactification.

% ======================================================================
\section{Summary Table and Status}
\label{sec:summary}
% ======================================================================

\begin{center}
\begin{tabular}{lllll}
  \toprule
  Candidate & Domain & $d_{\mathrm{eff}}$ & $N(\lambda)$ & Status \\
  \midrule
  1: $\Delta^{(4+1)}$ & $\mathbb{R}^4 \times S^1$ & 5 & $\lambda^{5/2}$ & [WRONG exponent] \\
  2: $\Delta^{(S^1)}$ & $S^1_\psi$ & 1 & $\lambda^{1/2}$ & [WRONG exponent] \\
  3: $\Delta^{(3)}_{\mathrm{Im}\,\mathbb{H}}$ & $\mathrm{Im}(\mathbb{H}) \cong \mathbb{R}^3$ & 3 & $\lambda^{3/2}$ & [CORRECT exponent — CANDIDATE] \\
  \bottomrule
\end{tabular}
\end{center}

\bigskip

\textbf{Key gap for Candidate 3:}

\begin{gap}[Spectral role of $\mathrm{Im}(\mathbb{H})$]
\label{gap:spectral_role}
  $\mathrm{Im}(\mathbb{H})$ is a 3-dimensional \emph{linear space} [PROVED].
  However, the statement that mode counting in UBT is governed by a
  Laplacian on $\mathrm{Im}(\mathbb{H})$ as a \emph{spectral manifold}
  is not yet proved.
  Specifically: the UBT vacuum polarisation is a spacetime integral
  (over $\mathbb{R}^4 \times S^1_\psi$), and the three factors of
  $\mathrm{Im}(\mathbb{H})$ enter as \emph{multiplicity factors},
  not as directions of a spectral manifold.
  The exponent $3/2$ in $N(\lambda) \sim \lambda^{3/2}$ requires
  $\mathrm{Im}(\mathbb{H})$ to contribute to $d_{\mathrm{eff}}$
  multiplicatively (as multiplicity = 3) \emph{and} with the correct
  normalisation to give the exponent $3/2 = 3 \times \frac{1}{2}$
  from the Gaussian path-integral identity.
  This identification must be made explicit.
\end{gap}

\bigskip

\textbf{Recommended forward path:}
Treat Candidate~3 as the working hypothesis in
\texttt{heat\_kernel\_theta.tex} (Step~3) and
\texttt{weyl\_counting\_theta.tex} (Step~4),
while being explicit about where the gap in Gap~\ref{gap:spectral_role}
enters.

\ifdefined\INCLUDEMODE\else\end{document}\fi
